\documentclass{article}

\usepackage[a4paper]{geometry}
\usepackage[english]{babel}
\usepackage{parskip}

\usepackage{amsmath, amssymb}
\usepackage{graphicx}
\usepackage{subcaption}
\usepackage{float}
\usepackage[hidelinks]{hyperref}

\usepackage{listings}
\usepackage{xcolor}

\definecolor{codegreen}{rgb}{0,0.6,0}
\definecolor{codegray}{rgb}{0.5,0.5,0.5}
\definecolor{codepurple}{rgb}{0.58,0,0.82}
\definecolor{backcolour}{rgb}{0.95,0.95,0.92}

\lstdefinestyle{mystyle}{
  backgroundcolor=\color{backcolour},
  commentstyle=\color{codegreen},
  keywordstyle=\color{magenta},
  numberstyle=\tiny\color{codegray},
  stringstyle=\color{codepurple},
  basicstyle=\ttfamily\footnotesize,
  breaklines=true,
  captionpos=b,
  keepspaces=true,
  numbers=left,
  numbersep=5pt,
  showstringspaces=false,
  tabsize=2
}
\lstset{style=mystyle}

\newcommand{\CodeDir}{Code}
\newcommand{\PlotDir}{Plots}

\title{NURa Hand-in 1}
\author{Silvan Toet (s1827332)}
\date{\today}

\begin{document}
\maketitle

\section{Poisson Distribution}
In this question, we are asked to compute the Poisson distribution for integer k:
\begin{equation}
  P_\lambda(k) = \frac{\lambda^k e^{-\lambda}}{k!}
\end{equation}
given a positive mean $\lambda$. We are asked to compute $P_\lambda(k)$ for the values in \ref{tab:poisson}.
We are restricted to 32 bit numbers. This means we will get overflow errors for most combinations of $k$ and $\lambda$ due to the rapid growth of $k!$ and $\lambda^{k}$ and underflow due to $e^{-lambda}$.
To prevent this we rewrite the equation in log10 form:
\begin{align}
  \\
  log_{10}(P_\lambda(k)) & = log_{10}(\lambda^k) + log_{10}(e^{-\lambda}) - log_{10}(k!)                    \\
                         & = k\ log_{10}(\lambda) - \lambda\ log_{10}(e) - log_{10}(\prod_{i=1}^{k} i)      \\
                         & = k\ log_{10}(\lambda) - \lambda\ log_{10}(e) - \sum_{i=1}^{k}i\ log_{10}        \\
  -> P_\lambda(k)        & = 10^{(k\ log_{10}(\lambda) - \lambda\ log_{10}(e) - \sum_{i=1}^{k}i\ log_{10})}
\end{align}
We choose log10 over the natural logarithm as this leads to a larger 'compression' of the large numbers, preventing overflow for a wider input range.


% Add the full code below.
The full code used for this question:
\lstinputlisting[language=Python]{\CodeDir/Poisson_code.txt}

% Show your results.
The results of $P_\lambda(k)$ for selected values of $k$ and $\lambda$ are shown in Table~\ref{tab:poisson}.

\begin{table}[H]
  \centering
  \caption{Poisson probability distribution for selected $k$ and $\lambda$.}
  \label{tab:poisson}
  \begin{tabular}{|c|c|c|}
    \hline
    $\lambda$ & $k$ & $P_\lambda(k)$ \\
    \hline
    \input{Poisson_output.txt}       \\
    \hline
  \end{tabular}
\end{table}

\section{Vandermonde Matrix and Interpolation}
For question 2 we are asked to solve the equation $\boldsymbol{Vc = y}$ for c.
This gives us the coeffients of the unique polynomial that passes through all points:
\begin{align}
  y_i = \sum_{j=0}^{N-1} c_j x_i^j (i = 0, ..., N-1)
\end{align}

First we construct the Vandermonde matrix V using x values in the supplied datafile.
Then in 2a we first create the LU decomposition of V with Crout's algorithm.
Next we solve for the supplied y values to find the coefficients c of the unique polynomial.
In 2b we use polynomial interpolation using Neville's algorithm to compare our results.
In 2c we use iterative improvement of the LU decomposition to get a better estimate.
In 2d we time the different methods. After printing all our code we will discuss our results.
\subsection{(2a) LU decomposition}

\begin{figure}[H]
  \centering
  \includegraphics[width=0.80\textwidth]{\PlotDir/vandermonde_sol_2a.pdf}
  \caption{Polynomial fit evaluated using LU decomposition. Top: data points and interpolated curve. Bottom: absolute error at the data points on a log scale.}
  \label{fig:vandermonde_2a}
\end{figure}

\subsection{(2b) Neville's algorithm}
\begin{figure}[H]
  \centering
  \includegraphics[width=0.80\textwidth]{\PlotDir/vandermonde_sol_2b.pdf}
  \caption{Interpolation using Neville's algorithm. Top: data points and interpolated curve. Bottom: absolute error at the data points on a log scale.}
  \label{fig:vandermonde_2b}
\end{figure}

\subsection{(2c) Improving the LU decomposition}
\begin{figure}[H]
  \centering
  \includegraphics[width=0.80\textwidth]{\PlotDir/vandermonde_sol_2c.pdf}
  \caption{LU-based solution with iterative refinement (showing iterations 0, 1 and 10). Top: interpolated curves. Bottom: absolute error at the data points on a log scale.}
  \label{fig:vandermonde_2c}
\end{figure}

\subsection{(2d) Timing}

Timing results (average per run):
\begin{itemize}
  \input{Execution_times.txt}
\end{itemize}

\subsection{Code for Question 2}
The following code was used for parts (2a)--(2d):
\lstinputlisting[language=Python]{\CodeDir/Vandermonde_all_code.txt}

\subsection{Conclusions}
As we can see from the plots in this section, Neville's algorithm is most accurate.
The LU decomposition method starts to oscillate more and more as we increase our x values.
This is likely due to the fact that the Vandermonde matrix becomes more and more ill-conditioned as we increase the x values, leading to larger numerical errors in the solution.
I think this problem is amplified due to my method not implementing partial pivoting, which would help to reduce the numerical errors in the solution.
The iterative improvement of the LU decomposition does improve the accuracy, but it still does not reach the accuracy of Neville's algorithm.
This is likely because the initial solution obtained from the LU decomposition is already quite inaccurate, and the iterative improvement can only do so much to correct it.

However, in terms of speed, LU decomposition is much faster than Neville's algorithm in this case, roughly by a factor of 50.
We also see that 10 extra iterations to improve the LU decomposition solution only doubles the execution time.
Most of the execution time is spent on constructing the LU decomposition itself.
\end{document}

